% ==============================================================================
% chapter2.tex — Глава 2. Техническая реализация
% ==============================================================================

\chapter{{\fontsize{14pt}{21pt}\selectfont\MakeUppercase{Техническая реализация}}}
\label{ch:chapter2}

\section{Проектирование архитектуры системы}
\label{sec:architecture}

На основании требований, сформированных в~первой главе, была спроектирована архитектура системы OneVoice. Ключевым архитектурным решением стало применение микросервисного подхода, обусловленного следующими факторами:

\begin{enumerate}
  \item необходимость независимого масштабирования отдельных компонентов (платформенных агентов);
  \item изоляция сбоев~--- отказ одного агента не влияет на работу остальных;
  \item возможность использования различных технологий для разных компонентов (Go для серверной части, Node.js для клиентского приложения, Playwright для агента браузерной автоматизации);
  \item упрощение добавления новых платформ без модификации существующего кода.
\end{enumerate}

Система состоит из шести микросервисов и разделяемой библиотеки. Общая архитектура представлена на диаграмме компонентов (рисунок~\ref{fig:component-diagram}).

\begin{figure}[H]
  \centering
  \includegraphics[width=\textwidth]{component-diagram.png}
  \caption{Диаграмма компонентов системы OneVoice}
  \label{fig:component-diagram}
\end{figure}

Состав и назначение микросервисов приведены в~таблице~\ref{tab:services}.

\begin{table}[H]
  \caption{Микросервисы системы OneVoice}
  \label{tab:services}
  \small
  \begin{tabularx}{\textwidth}{|p{3cm}|X|p{3.5cm}|}
    \hline
    \textbf{Модуль} & \textbf{Назначение} & \textbf{Технологии} \\
    \hline
    API-сервис & REST~API, аутентификация, операции над бизнес-данными (CRUD), проксирование чата & Go, Chi, PostgreSQL, MongoDB \\
    \hline
    Оркестратор & Агентный цикл LLM, диспетчеризация инструментов через NATS & Go, Chi, NATS \\
    \hline
    Веб-приложение & Панель управления (панель управления) & Next.js~14, React, TypeScript \\
    \hline
    Telegram-агент & Интеграция с~Telegram Bot API & Go, Telegram Bot API \\
    \hline
    VK-агент & Интеграция с~VK API & Go, VK API \\
    \hline
    Yandex.Business-агент & Интеграция с~Яндекс.Бизнес через браузерную автоматизацию & Go, Playwright \\
    \hline
    Общая библиотека (pkg) & Доменные модели, маршрутизатор LLM-запросов, программный каркас протокола A2A & Go \\
    \hline
  \end{tabularx}
\end{table}

\subsection{Протоколы межсервисного взаимодействия}
\label{subsec:protocols-impl}

В~системе используются три протокола взаимодействия, характеристики которых приведены в~таблице~\ref{tab:protocols}.

\begin{table}[H]
  \caption{Протоколы взаимодействия между компонентами системы}
  \label{tab:protocols}
  \small
  \begin{tabularx}{\textwidth}{|p{2.5cm}|p{3.5cm}|p{2.5cm}|X|}
    \hline
    \textbf{Протокол} & \textbf{Направление} & \textbf{Паттерн} & \textbf{Назначение} \\
    \hline
    REST (HTTP/JSON) & Веб-приложение $\rightarrow$ API-сервис & Запрос/ответ & CRUD-операции, аутентификация \\
    \hline
    SSE & API-сервис $\rightarrow$ Веб-приложение & Однонаправленный поток & Потоковая передача ответов LLM и статусов инструментов \\
    \hline
    NATS & Оркестратор $\rightarrow$ Агенты & Запрос/ответ & Диспетчеризация задач платформенным агентам \\
    \hline
  \end{tabularx}
\end{table}

Взаимодействие между компонентами при основных сценариях использования представлено на диаграмме последовательности (рисунок~\ref{fig:sequence-diagram}).

\begin{figure}[H]
  \centering
  \includegraphics[width=\textwidth]{sequence-diagram.png}
  \caption{Диаграмма последовательности: обновление информации и сбор отзывов}
  \label{fig:sequence-diagram}
\end{figure}

\textbf{Сценарий~1: обновление информации о~бизнесе.} Пользователь вводит в~чат запрос на обновление описания бизнеса на всех площадках. Через API-сервис сообщение проксируется к~оркестратору, который формирует контекст (профиль бизнеса, список активных интеграций) и отправляет запрос к~LLM с~описанием доступных инструментов. LLM последовательно определяет необходимые инструменты (\texttt{vk\_\_update\_community\_info}, \texttt{yandex\_business\_\_update\_info}) и возвращает вызовы инструментов. Оркестратор маршрутизирует каждую задачу через NATS соответствующему агенту, который выполняет действие на платформе. После получения результатов всех вызовов LLM формирует текстовый ответ.

\textbf{Сценарий~2: сбор отзывов.} Пользователь запрашивает новые отзывы. LLM вызывает инструмент \texttt{yandex\_business\_\_get\_reviews}, RPA-агент собирает отзывы с~платформы и возвращает их оркестратору. LLM анализирует полученные данные и формирует сводку для пользователя. Весь процесс транслируется клиенту в~реальном времени через SSE.

\subsection{Диаграмма развёртывания}
\label{subsec:deployment}

Все компоненты системы контейнеризованы с~использованием Docker и управляются через Docker~Compose. Диаграмма развёртывания представлена на рисунке~\ref{fig:deployment-diagram}.

\begin{figure}[H]
  \centering
  \includegraphics[width=\textwidth]{deployment-diagram.png}
  \caption{Диаграмма развёртывания системы OneVoice}
  \label{fig:deployment-diagram}
\end{figure}

Инфраструктурные компоненты включают: PostgreSQL для хранения реляционных данных, MongoDB для хранения документов и NATS в~качестве брокера сообщений.

\subsection{Маршрутизация запросов к LLM-провайдерам}
\label{subsec:llm-router}

Для взаимодействия с~большими языковыми моделями реализован LLM-маршрутизатор, поддерживающий несколько провайдеров: OpenRouter (основной, мультимодельный), OpenAI (прямое подключение), Anthropic (прямое подключение) и Self-Hosted (OpenAI-совместимые конечные точки). Маршрутизатор обеспечивает единый интерфейс для оркестратора независимо от используемого провайдера, что позволяет переключаться между моделями без изменения логики приложения.

\subsection{Протокол Agent-to-Agent (A2A)}
\label{subsec:a2a-impl}

Для стандартизации взаимодействия между оркестратором и платформенными агентами разработан протокол A2A, определяющий две структуры данных:

\begin{enumerate}
  \item \texttt{ToolRequest}~--- запрос от оркестратора к~агенту, содержащий идентификатор задачи, имя инструмента, аргументы и идентификатор бизнеса;
  \item \texttt{ToolResponse}~--- ответ агента, содержащий статус выполнения, результат или описание ошибки.
\end{enumerate}

Именование инструментов следует конвенции \texttt{\{платформа\}\_\_\{действие\}}, например: \texttt{telegram\_\_send\_channel\_post}, \texttt{vk\_\_publish\_wall\_post}, \texttt{yandex\_business\_\_update\_info}. Такой подход позволяет оркестратору автоматически определять целевого агента по префиксу имени инструмента и маршрутизировать задачу через NATS на соответствующий субъект (\texttt{tasks.telegram}, \texttt{tasks.vk}, \texttt{tasks.yandex\_business}).

\subsection{Алгоритм агентного цикла оркестратора}
\label{subsec:agent-loop}

Центральным элементом системы является агентный цикл оркестратора, реализующий итеративное взаимодействие с~LLM и платформенными агентами. Алгоритм представлен на рисунке~\ref{fig:activity-agent-loop}.

\begin{figure}[H]
  \centering
  \includegraphics[width=0.85\textwidth]{activity-agent-loop.png}
  \caption{Диаграмма деятельности: агентный цикл оркестратора}
  \label{fig:activity-agent-loop}
\end{figure}

Алгоритм работает следующим образом:

\begin{enumerate}
  \item При получении сообщения пользователя оркестратор формирует контекст, включающий профиль бизнеса, список активных интеграций и историю диалога.
  \item Контекст дополняется системной инструкцией (system prompt) с~описанием доступных инструментов и отправляется к~LLM-провайдеру.
  \item LLM анализирует запрос и возвращает либо текстовый ответ, либо набор вызовов инструментов (tool calls).
  \item Если ответ содержит вызовы инструментов, оркестратор извлекает имя инструмента и аргументы, определяет целевого агента по префиксу имени и отправляет запрос \texttt{ToolRequest} через NATS.
  \item Результат выполнения (\texttt{ToolResponse}) добавляется в~контекст диалога, и цикл повторяется.
  \item Цикл завершается, когда LLM возвращает текстовый ответ без вызовов инструментов или при достижении максимального числа итераций (\texttt{MAX\_ITERATIONS~=~10}).
\end{enumerate}

Каждое промежуточное состояние (вызов инструмента, результат, текст) транслируется клиенту в~реальном времени через SSE-события, обеспечивая прозрачность процесса для пользователя.

\subsection{Алгоритм ротации JWT-токенов}
\label{subsec:jwt-rotation}

Для обеспечения безопасной аутентификации в~системе реализован механизм ротации JWT-токенов. Алгоритм представлен на рисунке~\ref{fig:activity-jwt-rotation}.

\begin{figure}[H]
  \centering
  \includegraphics[width=0.85\textwidth]{activity-jwt-rotation.png}
  \caption{Диаграмма деятельности: ротация JWT-токенов}
  \label{fig:activity-jwt-rotation}
\end{figure}

Механизм аутентификации основан на паре токенов:

\begin{enumerate}
  \item \textbf{Access-токен}~--- JWT со сроком действия 15~минут, содержащий идентификатор пользователя и роль. Используется для авторизации каждого API-запроса.
  \item \textbf{Refresh-токен}~--- случайная строка со сроком действия 30~дней. Хранится в~базе данных в~виде хеша (SHA-256) и используется для обновления пары токенов.
\end{enumerate}

При обновлении токенов применяется ротация: старый refresh-токен удаляется из базы данных, а~новый сохраняется. Это предотвращает повторное использование скомпрометированных токенов. При выходе пользователя из системы refresh-токен отзывается, что делает невозможным дальнейшее обновление сессии.


\section{Проектирование модели данных}
\label{sec:data-model}

При проектировании модели данных использован подход гетерогенного хранения данных (polyglot persistence): реляционная СУБД PostgreSQL применяется для структурированных данных с~ACID-гарантиями (пользователи, бизнесы, интеграции), а~документоориентированная СУБД MongoDB~--- для данных с~гибкой схемой (диалоги, сообщения, публикации, отзывы, задачи агентов).

\subsection{Инфологическая модель (ER-диаграмма)}
\label{subsec:er-model}

Инфологическая модель реляционной части системы представлена на рисунке~\ref{fig:er-diagram-ch2}.

\begin{figure}[H]
  \centering
  \includegraphics[width=0.95\textwidth]{er-diagram.png}
  \caption{ER-диаграмма реляционной модели данных}
  \label{fig:er-diagram-ch2}
\end{figure}

Модель включает следующие сущности:

\begin{enumerate}
  \item \textbf{Пользователь (User)}~--- учётная запись с~аутентификационными данными и ролью в~системе;
  \item \textbf{Бизнес (Business)}~--- профиль бизнеса, связанный с~пользователем отношением <<один ко многим>>;
  \item \textbf{Интеграция (Integration)}~--- подключение к~внешней платформе с~зашифрованными токенами доступа (AES-256-GCM);
  \item \textbf{Расписание (BusinessSchedule)}~--- график работы бизнеса с~поддержкой особых дат;
  \item \textbf{Токен обновления (RefreshToken)}~--- хешированный токен для ротации JWT.
\end{enumerate}

\subsection{Даталогическая модель PostgreSQL}
\label{subsec:pg-schema}

Структура основных таблиц реляционной базы данных приведена в~таблицах~\ref{tab:users}--\ref{tab:refresh-tokens}.

\begin{table}[H]
  \caption{Структура таблицы users}
  \label{tab:users}
  \small
  \begin{tabularx}{\textwidth}{|p{3cm}|p{3.5cm}|X|}
    \hline
    \textbf{Поле} & \textbf{Тип} & \textbf{Описание} \\
    \hline
    id & UUID & Первичный ключ, генерируется автоматически \\
    \hline
    email & TEXT (UNIQUE) & Электронная почта пользователя \\
    \hline
    password\_hash & TEXT & Хеш пароля (bcrypt) \\
    \hline
    role & TEXT & Роль: owner, admin, member \\
    \hline
    created\_at & TIMESTAMPTZ & Дата создания записи \\
    \hline
    updated\_at & TIMESTAMPTZ & Дата последнего обновления \\
    \hline
  \end{tabularx}
\end{table}

\begin{table}[H]
  \caption{Структура таблицы businesses}
  \label{tab:businesses}
  \small
  \begin{tabularx}{\textwidth}{|p{3cm}|p{3cm}|X|}
    \hline
    \textbf{Поле} & \textbf{Тип} & \textbf{Описание} \\
    \hline
    id & UUID & Первичный ключ \\
    \hline
    user\_id & UUID (FK) & Ссылка на владельца (users.id) \\
    \hline
    name & TEXT & Название бизнеса \\
    \hline
    category & TEXT & Категория деятельности \\
    \hline
    phone & TEXT & Контактный телефон \\
    \hline
    website & TEXT & Адрес веб-сайта \\
    \hline
    address & TEXT & Физический адрес \\
    \hline
    description & TEXT & Краткое описание \\
    \hline
    logo\_url & TEXT & URL логотипа \\
    \hline
    settings & JSONB & Дополнительные настройки \\
    \hline
  \end{tabularx}
\end{table}

\begin{table}[H]
  \caption{Структура таблицы integrations}
  \label{tab:integrations}
  \small
  \begin{tabularx}{\textwidth}{|p{4cm}|p{3cm}|X|}
    \hline
    \textbf{Поле} & \textbf{Тип} & \textbf{Описание} \\
    \hline
    id & UUID & Первичный ключ \\
    \hline
    business\_id & UUID (FK) & Ссылка на бизнес \\
    \hline
    platform & TEXT & Платформа: telegram, vk, yandex\_business \\
    \hline
    status & TEXT & Статус: pending, active, error \\
    \hline
    encrypted\_access\_token & BYTEA & Токен доступа (AES-256-GCM) \\
    \hline
    encrypted\_refresh\_token & BYTEA & Токен обновления (AES-256-GCM) \\
    \hline
    external\_id & TEXT & Внешний идентификатор (ID канала, группы) \\
    \hline
    metadata & JSONB & Метаданные подключения \\
    \hline
    token\_expires\_at & TIMESTAMPTZ & Срок действия токена \\
    \hline
  \end{tabularx}
\end{table}

Ограничение уникальности \texttt{UNIQUE (business\_id, platform, external\_id)} позволяет подключать несколько каналов одной платформы к~одному бизнесу, что важно для работы с~несколькими Telegram-каналами или VK-сообществами.

\begin{table}[H]
  \caption{Структура таблицы business\_schedules}
  \label{tab:schedules}
  \small
  \begin{tabularx}{\textwidth}{|p{3cm}|p{3cm}|X|}
    \hline
    \textbf{Поле} & \textbf{Тип} & \textbf{Описание} \\
    \hline
    id & UUID & Первичный ключ \\
    \hline
    business\_id & UUID (FK) & Ссылка на бизнес (businesses.id) \\
    \hline
    day\_of\_week & INTEGER & День недели (0~--- воскресенье, 6~--- суббота) \\
    \hline
    is\_open & BOOLEAN & Признак рабочего дня \\
    \hline
    open\_time & TEXT & Время открытия (формат HH:MM) \\
    \hline
    close\_time & TEXT & Время закрытия (формат HH:MM) \\
    \hline
    special\_date & DATE & Особая дата (если задана, запись описывает исключение из обычного расписания) \\
    \hline
  \end{tabularx}
\end{table}

\begin{table}[H]
  \caption{Структура таблицы refresh\_tokens}
  \label{tab:refresh-tokens}
  \small
  \begin{tabularx}{\textwidth}{|p{3cm}|p{3cm}|X|}
    \hline
    \textbf{Поле} & \textbf{Тип} & \textbf{Описание} \\
    \hline
    id & UUID & Первичный ключ \\
    \hline
    user\_id & UUID (FK) & Ссылка на пользователя (users.id) \\
    \hline
    token\_hash & TEXT & Хеш refresh-токена (SHA-256) \\
    \hline
    expires\_at & TIMESTAMPTZ & Срок действия токена \\
    \hline
    created\_at & TIMESTAMPTZ & Дата создания записи \\
    \hline
  \end{tabularx}
\end{table}

\subsection{Модель документов MongoDB}
\label{subsec:mongo-schema}

Для хранения данных с~гибкой структурой используется MongoDB. Схема коллекций представлена на рисунке~\ref{fig:mongo-schema}.

\begin{figure}[H]
  \centering
  \includegraphics[width=0.85\textwidth]{mongo-schema.png}
  \caption{Схема коллекций MongoDB}
  \label{fig:mongo-schema}
\end{figure}

Основные коллекции:

\begin{enumerate}
  \item \textbf{conversations}~--- диалоги пользователя с~ИИ-ассистентом;
  \item \textbf{messages}~--- сообщения в~диалогах с~поддержкой вложений, вызовов инструментов и их результатов;
  \item \textbf{posts}~--- публикации на платформах с~результатами по каждой платформе;
  \item \textbf{reviews}~--- отзывы клиентов с~различных платформ;
  \item \textbf{agent\_tasks}~--- журнал задач, выполняемых платформенными агентами.
\end{enumerate}

Выбор MongoDB для этих коллекций обусловлен необходимостью хранения вложенных объектов переменной структуры (массивы вызовов инструментов в~сообщениях, результаты публикаций по платформам) без необходимости нормализации.

\subsection{Типовые запросы к базам данных}
\label{subsec:db-queries}

В~таблице~\ref{tab:queries} приведены типовые запросы, используемые в~системе.

\begin{table}[H]
  \caption{Типовые запросы к базам данных}
  \label{tab:queries}
  \small
  \begin{tabularx}{\textwidth}{|p{3.5cm}|p{1.5cm}|X|}
    \hline
    \textbf{Операция} & \textbf{СУБД} & \textbf{Описание} \\
    \hline
    Аутентификация & PG & Поиск пользователя по email, проверка хеша пароля \\
    \hline
    Получение токена интеграции & PG & Поиск активной интеграции по business\_id и platform с~расшифровкой токена \\
    \hline
    Загрузка истории чата & Mongo & Получение сообщений по conversation\_id с~вложенными tool\_calls и tool\_results \\
    \hline
    Фильтрация публикаций & Mongo & Поиск по business\_id с~фильтрами по статусу и платформе, сортировка по дате \\
    \hline
    Фильтрация отзывов & Mongo & Поиск по business\_id с~фильтрами по платформе и статусу ответа \\
    \hline
    Журнал задач агентов & Mongo & Поиск по business\_id с~фильтрами по статусу и платформе \\
    \hline
  \end{tabularx}
\end{table}


\section{Проектирование пользовательского интерфейса}
\label{sec:ui-design}

Пользовательский интерфейс системы OneVoice реализован как одностраничное веб-приложение (SPA) на базе Next.js~14 с~использованием компонентной библиотеки shadcn/ui. При проектировании интерфейса учитывались следующие принципы:

\begin{enumerate}
  \item \textbf{Единообразие}~--- все страницы используют общий макет с~боковой навигационной панелью;
  \item \textbf{Адаптивность}~--- интерфейс корректно отображается на настольных компьютерах, планшетах и мобильных устройствах;
  \item \textbf{Обратная связь}~--- действия пользователя сопровождаются визуальными индикаторами (статусы подключений, уведомления, анимации загрузки);
  \item \textbf{Минимализм}~--- интерфейс не перегружен элементами, основной акцент на функциональности.
\end{enumerate}

Выбор архитектуры одностраничного приложения (SPA) обусловлен спецификой основного сценария использования~--- диалога с~ИИ-ассистентом в~реальном времени. SPA позволяет поддерживать постоянное SSE-соединение для потоковой передачи ответов LLM без перезагрузки страницы, а~также обеспечивает мгновенные переходы между разделами панели управления с~сохранением состояния чата.

В~качестве основного цвета выбран оттенок Indigo (\#6366F1), обеспечивающий профессиональное восприятие интерфейса и высокий контраст с~белым фоном (коэффициент контрастности 4.6:1 по WCAG~2.1~AA). Данный цвет является стандартным для библиотеки shadcn/ui, что гарантирует визуальную согласованность всех компонентов.

Шрифтовое оформление построено на гарнитуре Inter~--- свободном шрифте, оптимизированном для экранного отображения и поддерживающем кириллицу. Inter входит в~стандартный набор системных шрифтов Next.js, что исключает необходимость загрузки дополнительных веб-шрифтов и ускоряет начальную отрисовку страницы.

\subsection{Структура навигации}
\label{subsec:navigation}

Система разделена на две зоны: публичную (аутентификация) и защищённую (основной панель управления). Экранные формы системы перечислены в~таблице~\ref{tab:screens}.

\begin{table}[H]
  \caption{Экранные формы системы OneVoice}
  \label{tab:screens}
  \small
  \begin{tabularx}{\textwidth}{|p{2.5cm}|p{3cm}|X|}
    \hline
    \textbf{Страница} & \textbf{Назначение} & \textbf{Ключевые элементы} \\
    \hline
    Вход & Аутентификация & Форма: email, пароль, кнопка входа \\
    \hline
    Регистрация & Создание аккаунта & Форма: имя, email, пароль, подтверждение \\
    \hline
    Чат & Диалог с~ИИ-ассистентом & Список диалогов, окно сообщений, панель инструментов, быстрые действия \\
    \hline
    Интеграции & Подключение платформ & Карточки платформ, модальные окна подключения, статусы \\
    \hline
    Бизнес & Профиль и расписание & Форма профиля, редактор расписания по дням, особые даты \\
    \hline
    Отзывы & Управление отзывами & Таблица с~фильтрами, встроенный диалог ответа \\
    \hline
    Посты & Журнал публикаций & Таблица с~фильтрами, раскрываемые строки с~результатами по платформам \\
    \hline
    Задачи & Журнал задач агентов & Таблица с~фильтрами по статусу и платформе \\
    \hline
    Настройки & Параметры аккаунта & Информация о~пользователе, смена пароля \\
    \hline
  \end{tabularx}
\end{table}

Навигация по защищённой зоне осуществляется через боковую панель шириной 240~пикселей, содержащую ссылки на все разделы и индикаторы подключения платформ в~нижней части. На мобильных устройствах панель скрывается и вызывается жестом или кнопкой.

\subsection{Цветовое решение и типографика}
\label{subsec:colors}

Цветовая палитра и параметры типографики приведены в~таблице~\ref{tab:design-tokens}.

\begin{table}[H]
  \caption{Цветовая палитра и типографика интерфейса}
  \label{tab:design-tokens}
  \small
  \begin{tabularx}{\textwidth}{|p{4cm}|p{3.5cm}|X|}
    \hline
    \textbf{Параметр} & \textbf{Значение} & \textbf{Применение} \\
    \hline
    Основной цвет (Primary) & \#6366F1 (Indigo) & Кнопки, активные элементы навигации, ссылки \\
    \hline
    Фон боковой панели & \#1E293B (Slate~800) & Навигационная панель \\
    \hline
    Фон контентной области & \#F8FAFC (Slate~50) & Основная область страницы \\
    \hline
    Цвет успеха & \#22C55E (Green) & Статусы <<Подключено>>, <<Опубликовано>> \\
    \hline
    Цвет ошибки & \#EF4444 (Red) & Ошибки, статусы <<Отключено>> \\
    \hline
    Шрифт & Inter (системный) & Основной текст, заголовки \\
    \hline
    Размер базового текста & 14px & Основной текст интерфейса \\
    \hline
  \end{tabularx}
\end{table}

\subsection{Экранные формы}
\label{subsec:screens-detail}

\subsubsection{Аутентификация}

Страницы входа и регистрации выполнены в~минималистичном стиле с~центрированной формой на светлом фоне (рисунки~\ref{fig:ui-login},~\ref{fig:ui-register}).

\begin{figure}[H]
  \centering
  \includegraphics[width=0.6\textwidth]{ui-login.png}
  \caption{Экранная форма авторизации}
  \label{fig:ui-login}
\end{figure}

\begin{figure}[H]
  \centering
  \includegraphics[width=0.6\textwidth]{ui-register.png}
  \caption{Экранная форма регистрации}
  \label{fig:ui-register}
\end{figure}

\subsubsection{Профиль бизнеса и расписание}

Страница профиля бизнеса содержит форму с~основной информацией (название, категория, телефон, сайт, адрес, описание) и редактор расписания работы по дням недели с~поддержкой выходных дней и особых дат (рисунки~\ref{fig:ui-business},~\ref{fig:ui-schedule}).

\begin{figure}[H]
  \centering
  \includegraphics[width=\textwidth]{ui-business.png}
  \caption{Экранная форма профиля бизнеса с~расписанием}
  \label{fig:ui-business}
\end{figure}

\begin{figure}[H]
  \centering
  \includegraphics[width=\textwidth]{ui-schedule.png}
  \caption{Редактор расписания работы}
  \label{fig:ui-schedule}
\end{figure}

\subsubsection{Управление интеграциями}

Страница интеграций отображает карточки доступных платформ (Telegram, ВКонтакте, Яндекс.Бизнес) с~индикаторами статуса подключения. Для каждой платформы предусмотрено модальное окно подключения с~пошаговой процедурой (рисунок~\ref{fig:ui-integrations}).

\begin{figure}[H]
  \centering
  \includegraphics[width=\textwidth]{ui-integrations.png}
  \caption{Экранная форма управления интеграциями}
  \label{fig:ui-integrations}
\end{figure}

\subsubsection{Журнал публикаций}

Страница публикаций содержит таблицу с~фильтрацией по платформе и статусу (все, опубликованы, запланированы, ошибки). Каждая строка раскрывается для отображения полного текста и результатов публикации по платформам (рисунки~\ref{fig:ui-posts},~\ref{fig:ui-posts-expanded}).

\begin{figure}[H]
  \centering
  \includegraphics[width=\textwidth]{ui-posts.png}
  \caption{Журнал публикаций}
  \label{fig:ui-posts}
\end{figure}

\begin{figure}[H]
  \centering
  \includegraphics[width=\textwidth]{ui-posts-expanded.png}
  \caption{Раскрытая запись публикации с~результатами по платформам}
  \label{fig:ui-posts-expanded}
\end{figure}

\subsubsection{Журнал задач агентов}

Страница задач предоставляет мониторинг операций, выполняемых платформенными агентами, с~фильтрацией по статусу (в~очереди, активные, завершены, ошибки) и платформе (рисунок~\ref{fig:ui-tasks}).

\begin{figure}[H]
  \centering
  \includegraphics[width=\textwidth]{ui-tasks.png}
  \caption{Журнал задач агентов}
  \label{fig:ui-tasks}
\end{figure}

\subsection{Управление состоянием}
\label{subsec:state-management}

Для управления состоянием клиентского приложения используется двухуровневый подход:

\begin{enumerate}
  \item \textbf{Zustand}~--- локальное состояние аутентификации (пользователь, токены, статус авторизации) с~сохранением состояния в~локальном хранилище браузера (localStorage);
  \item \textbf{TanStack React Query}~--- серверное состояние (диалоги, интеграции, бизнес-профиль, публикации, отзывы, задачи) с~автоматическим кешированием, инвалидацией и повторными запросами.
\end{enumerate}

Потоковая передача ответов LLM реализована через SSE (Server-Sent Events): клиент устанавливает соединение при отправке сообщения и обрабатывает события пяти типов: \texttt{text} (текст ответа), \texttt{tool\_call} (вызов инструмента), \texttt{tool\_result} (результат инструмента), \texttt{done} (завершение), \texttt{error} (ошибка).


\section{Средства реализации}
\label{sec:tech-impl}

Система OneVoice реализована на языках Go (серверная часть) и TypeScript (клиентское приложение). Общий объём исходного кода составляет порядка 26\,000~строк: 16\,000~строк Go-кода в~микросервисах, 4\,000~строк в~разделяемой библиотеке \texttt{pkg} и 6\,000~строк TypeScript/TSX в~веб-приложении.

Ключевые библиотеки и фреймворки, использованные при разработке, перечислены в~таблице~\ref{tab:libs}.

\begin{table}[H]
  \caption{Основные библиотеки и фреймворки}
  \label{tab:libs}
  \small
  \begin{tabularx}{\textwidth}{|p{2.5cm}|p{3.5cm}|X|}
    \hline
    \textbf{Платформа} & \textbf{Библиотека} & \textbf{Назначение} \\
    \hline
    Go & chi v5 & HTTP-маршрутизация и middleware \\
    \hline
    Go & pgx v5 & Драйвер PostgreSQL с~поддержкой пула соединений \\
    \hline
    Go & nats.go & Клиент NATS для межсервисной коммуникации \\
    \hline
    Go & playwright-go & Автоматизация браузера для RPA-агента \\
    \hline
    Go & golang-jwt v5 & Генерация и валидация JWT-токенов \\
    \hline
    TypeScript & Next.js~14 & SSR/SSG-фреймворк с~App Router \\
    \hline
    TypeScript & TanStack React Query & Управление серверным состоянием и кешированием \\
    \hline
    TypeScript & Zustand & Управление локальным состоянием (аутентификация) \\
    \hline
    TypeScript & Zod & Валидация данных форм и API-ответов \\
    \hline
    TypeScript & shadcn/ui & Компонентная библиотека на базе Radix~UI \\
    \hline
  \end{tabularx}
\end{table}

Системные требования для развёртывания: Docker Engine~20+, Docker Compose~v2, 4~ГБ оперативной памяти, Go~1.24 и Node.js~20 (для локальной разработки). Полный исходный код системы приведён в~приложении~А.


\section{Инсталляция и настройка}
\label{sec:installation}

Развёртывание системы OneVoice выполняется с~помощью Docker Compose и включает девять контейнеров: шесть микросервисов (API, оркестратор, веб-приложение, Telegram-агент, VK-агент, Yandex.Business-агент) и три инфраструктурных сервиса (PostgreSQL, MongoDB, NATS).

Конфигурация системы осуществляется через переменные окружения, определённые в~файле \texttt{.env}. Основные параметры включают строку подключения к~PostgreSQL (\texttt{DATABASE\_URL}), адрес MongoDB (\texttt{MONGODB\_URI}), адрес NATS (\texttt{NATS\_URL}), ключ шифрования токенов (\texttt{ENCRYPTION\_KEY}), секрет JWT (\texttt{JWT\_SECRET}), модель LLM (\texttt{LLM\_MODEL}) и API-ключи провайдеров.

Процедура развёртывания состоит из следующих шагов:

\begin{enumerate}
  \item Копирование файла \texttt{.env.example} в~\texttt{.env} и заполнение параметров.
  \item Запуск всех сервисов командой \texttt{make up}, которая выполняет \texttt{docker compose up -d}.
  \item Автоматическое применение миграций базы данных при первом запуске API-сервиса.
  \item Проверка работоспособности по адресу \texttt{http://localhost:3000}.
\end{enumerate}

Типовой сценарий начала работы пользователя включает: регистрацию учётной записи, создание профиля бизнеса с~заполнением информации и расписания, подключение интеграций с~платформами (ввод токена Telegram-бота, данных VK-сообщества, учётных данных Яндекс.Бизнес) и переход к~диалогу с~ИИ-ассистентом для управления цифровым присутствием.


\section{Выводы по главе 2}
\label{sec:ch2-conclusions}

В~данной главе выполнена техническая реализация системы OneVoice. Спроектирована микросервисная архитектура, включающая шесть компонентов (API-сервис, оркестратор, веб-приложение, три платформенных агента) и разделяемую библиотеку, взаимодействующих через REST, SSE и NATS.

Описаны два ключевых алгоритма: агентный цикл оркестратора, реализующий итеративное взаимодействие с~LLM и диспетчеризацию инструментов через протокол A2A, и алгоритм ротации JWT-токенов, обеспечивающий безопасную аутентификацию с~автоматическим обновлением сессии.

Модель данных построена на принципе гетерогенного хранения данных: реляционная СУБД PostgreSQL (5~таблиц: пользователи, бизнесы, интеграции, расписания, токены обновления) используется для структурированных данных с~ACID-гарантиями, а~документоориентированная MongoDB (5~коллекций: диалоги, сообщения, публикации, отзывы, задачи агентов)~--- для данных с~гибкой схемой.

Пользовательский интерфейс реализован как одностраничное веб-приложение на базе Next.js~14 и включает 10~экранных форм. Интерфейс спроектирован с~учётом принципов единообразия, адаптивности и минимализма, обеспечивая корректное отображение на настольных компьютерах, планшетах и мобильных устройствах.

Общий объём исходного кода составляет порядка 26\,000~строк. Развёртывание системы контейнеризовано с~помощью Docker~Compose (9~контейнеров), что обеспечивает воспроизводимость среды и упрощает установку.
